\section{Trace estimates uniform in the cutoff}
\label{sec:trace}

The endpoint comparison now produces a box of fixed physical size in
which the excited transfer trace is small at every sufficiently fine
cutoff.  Positivity then propagates this information to arbitrarily
large boxes.  This section proves both consequences needed below: a
uniform comparison between finite tori and the plane, and exponential
decorrelation at a fixed positive rate in physical units.  The transfer
argument follows \cite{OpenAI-O4}, \Rref{sec:finite-volume}; we give its
proof, including the rectangular and sheared periods used here.

Write $Z_\beta(n,w)$ for the nearest-neighbour partition function on the
torus with time period $n$ and spatial period $w$, using normalized area
measure at every spin.  All geometric periods below, including the
auxiliary half-periods, are at least four.  Define
\begin{equation}
 \Delta_\beta(n,w)
   =4\log Z_\beta(n,w)-\log Z_\beta(2n,2w).
 \label{eq:trace-doubling-definition}
\end{equation}
The coefficients cancel every contribution to $\log Z$ proportional to
the area.  In particular, the bulk scalars extracted by the exact
renormalization do not contribute to $\Delta_\beta$.

\subsection{A fixed physical box}

The gain in the preliminary length estimate and the contraction of
endpoint discrepancies have complementary roles.  The former makes a
reference box large enough to mix; the latter allows that box to be
compared with every finer cutoff.

\begin{proposition}[A uniform reference-box test]
\label{prop:trace-small-box}
Let $\mathcal U\subset(\mathbb Q_{>0})^2$ be finite.  Assume the
preliminary estimate of Section~\ref{sec:preliminary}, the shooting
asymptotics \eqref{eq:orig-12}, the endpoint discrepancy estimate
\eqref{eq:orig-15}, and the integrated comparison
\eqref{eq:orig-16}.  With the contraction slack $\upsilon>0$ chosen
sufficiently small, there exist an integer $K_0$ and a positive integer
$M_0$, independent of $N$, such that, for every $N\ge K_0$ and every
$(u_1,u_2)\in\mathcal U$,
\begin{equation}
 \Delta_{\beta_N}(n,w)\le 2^{-10},
 \qquad (n,w)=(u_1,u_2)M_0L^N.
 \label{eq:orig-24}
\end{equation}
The number $M_0$ may be required to exceed any prescribed fixed constant
and to be divisible by any prescribed fixed integer.  In particular,
all periods and tile counts needed in
Section~\ref{sec:comparison} may be taken to be even integers.
\end{proposition}

\begin{proof}
Let $\eta>0$ be the gain in \eqref{eq:orig-1}.  Choose
\[
 0<\frac{\upsilon}{2}<\chi<\min\left\{1,\frac{\eta}{2\pi}\right\}.
\]
Take $m_K$ to be a fixed common integer multiple of a dyadic integer,
rounded so that
\begin{equation}
 \log_L m_K=(1-\chi)K+O(1),\qquad
 2(1-\chi)<2-\upsilon,\qquad
 2-\chi>\frac{4\pi-\eta}{2\pi}.
 \label{eq:trace-scale-choice}
\end{equation}
The fixed multiple clears the denominators in $\mathcal U$ and imposes
the divisibility conditions of the proposition.

At depth $K$, Equation~\eqref{eq:orig-12} and
$\gamma=(\log L)/(2\pi)$ give
\[
 X(\beta_K)
   \le C_H K^{C_H} L^{(2-\eta/(2\pi))K}.
\]
The shorter side of a reference rectangle
$(u_1,u_2)m_KL^K$ is a fixed positive multiple of
$L^{(2-\chi)K}$.  Its ratio to $X(\beta_K)$ therefore grows
exponentially in $K$.  The preliminary doubling estimate consequently
implies
\[
 \max_{(u_1,u_2)\in\mathcal U}
 \left|\Delta_{\beta_K}(u_1m_KL^K,u_2m_KL^K)\right|
     \longrightarrow0.
\]
Indeed its exponentially decaying factor dominates every polynomial
prefactor in the microscopic periods and in $\beta_K$.

For $N\ge K$, run both cutoffs to the same terminal rectangle with
periods $(u_1,u_2)m_K$.  Its volume is
$v=u_1u_2m_K^2$.  Equations~\eqref{eq:orig-15} and
\eqref{eq:orig-16}, with the exponent discrepancy included, imply
that the logarithms of the two scalar-stripped partition functions
differ by at most $C_Hv\delta_K$ whenever $v\delta_K$ is sufficiently
small.  The estimate is uniform in $N\ge K$.  Apply it also to the
doubled rectangle, whose terminal volume is $4v$.  Cancellation of the
bulk scalars gives
\begin{align*}
 &\left|\Delta_{\beta_N}(u_1m_KL^N,u_2m_KL^N)
       -\Delta_{\beta_K}(u_1m_KL^K,u_2m_KL^K)\right|\\
 &\hspace{35mm}\le C_Hm_K^2\delta_K
       \longrightarrow0,
\end{align*}
because $\delta_K=L^{-(2-\upsilon)K}$ and
$2(1-\chi)<2-\upsilon$.  Constants may depend on the finite list
$\mathcal U$, but not on $N$ or $K$.  Choose one sufficiently large
$K_0$ and put $M_0=m_{K_0}$.
\end{proof}

\subsection{Positive transfer and repeated doubling}

We next isolate the spectral content of the partition-function test.
For a fixed spatial circumference $w$, a time slice is
$s=(s_1,\ldots,s_w)\in(S^2)^w$.  Put
$V_w(s)=\sum_{i=1}^w s_i\cdot s_{i+1}$, with periodic indices.
On $L^2((S^2)^w)$ define the transfer operator $K_w$ by the kernel
\begin{equation}
 K_w(s,s')=
 \exp\left\{\frac\beta2V_w(s)
       +\beta\sum_{i=1}^w s_i\cdot s_i'
       +\frac\beta2V_w(s')\right\}.
 \label{eq:trace-kernel}
\end{equation}
The coupling $\beta\ge0$ is fixed in this subsection and suppressed
from the notation.  All Hilbert spaces use product probability measure.

The kernel is continuous, symmetric, and strictly positive.  Its
positive semidefiniteness is a separate property: the expansion
\[
 e^{\beta s\cdot s'}
   =\sum_{j=0}^{\infty}\frac{\beta^j}{j!}
       \inner{s^{\otimes j}}{(s')^{\otimes j}}
\]
is a sum of positive semidefinite kernels.  Tensor products over sites
and multiplication on both sides by $e^{\beta V_w/2}$ preserve that
property.  The sum of their diagonal integrals is finite, so the same
expansion proves that $K_w$ is trace class.  Strict positivity of the
kernel gives a simple largest eigenvalue with a strictly positive
normalized eigenfunction $\Omega_w$.  Write
\[
 \lambda_1(w)>\lambda_2(w)\ge\lambda_3(w)\ge\cdots\ge0,
 \qquad
 Z_\beta(n,w)=\Tr K_w^n=\sum_i\lambda_i(w)^n.
\]
If only one eigenvalue is nonzero, all subsequent statements are read
with $\lambda_2=0$.  These observations also verify directly the
transfer hypotheses of \Rref{prop:criterion} for $S^2$ spins.

Define the two one-direction defects
\begin{align}
 p(n,w)&=2\log Z_\beta(n,w)-\log Z_\beta(2n,w),\notag\\
 q(n,w)&=2\log Z_\beta(n,w)-\log Z_\beta(n,2w).
 \label{eq:trace-one-direction}
\end{align}
Both are nonnegative, using transfer positivity in the respective
coordinate directions.  The excited trace relative to the largest
eigenvalue is
\[
 s_n(w)=\sum_{i\ge2}
       \left(\frac{\lambda_i(w)}{\lambda_1(w)}\right)^n.
\]

\begin{lemma}[Concentration of the transfer trace]
\label{lem:trace-concentration}
For every integer $n\ge4$,
$s_n(w)\le e^{p(n,w)}-1$.  If $p(n,w)\le1/4$, then
\begin{equation}
 \left(\frac{\lambda_2(w)}{\lambda_1(w)}\right)^n
       \le2p(n,w),\qquad
 p(2n,w)\le4p(n,w)^2.
 \label{eq:trace-squaring}
\end{equation}
The corresponding conclusions hold for $q$ with the coordinate
directions interchanged.
\end{lemma}

\begin{proof}
Set $t_i=\lambda_i(w)^n/\sum_j\lambda_j(w)^n$.  Then
\[
 e^{-p(n,w)}=\sum_i t_i^2\le t_1=(1+s_n(w))^{-1}.
\]
This gives the first assertion and, when $p\le1/4$, the bound
$s_n\le pe^p\le2p$.  Moreover $s_{2n}\le s_n^2$, and cancellation
of the largest eigenvalue yields
\[
 p(2n,w)=2\log(1+s_{2n})-\log(1+s_{4n})
     \le2s_n^2\le2p(n,w)^2e^{2p(n,w)}
     \le4p(n,w)^2.
\]
This is the proof of \Rref{lem:squaring}, with no dependence on the
dimension of the spin sphere.
\end{proof}

Small excited weight at one circumference alone would leave open the
appearance of low-energy states at larger circumferences.  The second
direction in $\Delta$ supplies the missing control.  Direct substitution
in \eqref{eq:trace-one-direction} gives
\begin{align}
 \Delta_\beta(n,w)
   &=2p(n,w)+q(2n,w)
     =2q(n,w)+p(n,2w),\notag\\
 p(n,2w)&=2p(n,w)-2q(n,w)+q(2n,w).
 \label{eq:trace-rectangle-identities}
\end{align}
In particular, $\Delta_\beta(n,w)\ge0$.

\begin{proposition}[Rectangular doubling criterion]
\label{prop:trace-doubling}
Let $\beta\ge0$ and let $n,w\ge4$ be integers.  If
$\Delta_\beta(n,w)\le2^{-10}$, then, for every $k\ge0$,
\begin{equation}
 \Delta_\beta(2^kn,2^kw)\le64^{-1}2^{-2^k}.
 \label{eq:orig-25}
\end{equation}
For a square with $n=w$, the transfer operators at circumferences
$2^kn$ also satisfy
\begin{equation}
 \frac{\lambda_2(2^kn)}{\lambda_1(2^kn)}
       \le e^{-(\log2)/n}.
 \label{eq:trace-width-gap}
\end{equation}
\end{proposition}

\begin{proof}
Write $\Delta=\Delta_\beta(n,w)$.  The two expressions in
\eqref{eq:trace-rectangle-identities} imply
$q(2n,w)\le\Delta$ and $p(n,2w)\le\Delta$.  Consequently
\[
 p(2n,2w)\le4\Delta^2.
\]
The interchanged rectangle identity then gives
\[
 q(4n,w)=2q(2n,w)-2p(2n,w)+p(2n,2w)
       \le2\Delta+4\Delta^2\le3\Delta.
\]
Squaring in the spatial direction gives
$q(4n,2w)\le36\Delta^2$, and hence
\[
 \Delta_\beta(2n,2w)
   =2p(2n,2w)+q(4n,2w)
   \le44\Delta^2\le64\Delta^2.
\]
Every application of Lemma~\ref{lem:trace-concentration} is permitted
by $3\cdot2^{-10}<1/4$, and the smallness condition is preserved.
Iteration yields
\[
 \Delta_\beta(2^kn,2^kw)
       \le64^{-1}(64\cdot2^{-10})^{2^k}
       \le64^{-1}2^{-2^k}.
\]
For a square, put $n_k=2^kn$.  The same lemma gives
\[
 \left(\frac{\lambda_2(n_k)}{\lambda_1(n_k)}\right)^{n_k}
       \le2p(n_k,n_k)\le\Delta_\beta(n_k,n_k).
\]
Taking the $n_k$-th root proves
\eqref{eq:trace-width-gap}.  Thus the square proof of
\Rref{prop:criterion} also proves the asserted rectangular recursion;
no equality of the original periods entered that recursion.
\end{proof}

\subsection{Local insertions and the volume limit}

To use the trace bound for correlations, one needs control of an
insertion that may occupy an entire time slab.  The following estimate
has no cost depending on the number of spins read by the insertion.
Let $F$ be a bounded local observable supported in a slab of $r$
transfer steps.  The kernel $K_{w,F}$ obtained by inserting $F$ in that
slab satisfies
\[
 |K_{w,F}(s,s')|\le\norm{F}_\infty K_w^r(s,s').
\]
With $T_w=K_w/\lambda_1(w)$ and
$A_F=K_{w,F}/\lambda_1(w)^r$, positivity of $K_w^r$ implies
$\norm{A_F}\le\norm{F}_\infty$: apply the pointwise inequality to
$|f|$ and use $\norm{T_w^r}=1$.  When $r=0$, $A_F$ is multiplication
by $F$.  The infinite-time cylinder expectation is therefore
\[
 \omega_w(F)=\inner{\Omega_w}{A_F\Omega_w}.
\]

Let $P_w$ be the orthogonal projection onto $\Omega_w$.  Since
$T_w-P_w$ is positive semidefinite,
$\Tr(T_w^j-P_w)=s_j(w)$.  For $r\le n/2$, splitting off $P_w$ in
\[
 \mu_{\beta;n,w}(F)
    =\frac{\Tr(A_FT_w^{n-r})}{\Tr T_w^n}
\]
gives
\begin{equation}
 |\mu_{\beta;n,w}(F)-\omega_w(F)|
      \le2\norm{F}_\infty s_{n/2}(w),
 \label{eq:trace-torus-cylinder}
\end{equation}
whenever $n$ is even.  Here $\mu_{\beta;n,w}$ denotes the probability
law on the rectangular torus.  The same calculation applies after
interchanging space and time.

We state the resulting volume comparison with the auxiliary aspect
ratios specified explicitly.  This makes the finite amount of
information required from Proposition~\ref{prop:trace-small-box}
independent of all later limits.  For complex observables the covariance
is sesquilinear:
$\Cov_\omega(F,G)=\omega(\overline F G)
 -\omega(\overline F)\omega(G)$.

\begin{theorem}[Uniform volume comparison and slab decorrelation]
\label{thm:trace-uniform}
Let $\mathcal A\subset(\mathbb Q_{>0})^2$ be a finite list containing
$(1,1)$ and $(1,25)$, and put
\[
 \mathcal U=\mathcal A
    \cup\{(u_1/2,u_2):(u_1,u_2)\in\mathcal A\}
    \cup\{(2u_1,u_2/2):(u_1,u_2)\in\mathcal A\}.
\]
Choose $K_0,M_0$ by Proposition~\ref{prop:trace-small-box} for
$\mathcal U$, with all displayed periods even and at least four.  Let
$\omega_{\beta_N}$ be the periodic infinite-plane state supplied by
Section~\ref{sec:preliminary}.  There are constants $C,c>0$, independent
of $N\ge K_0$ and $k\ge0$, with the following properties.

For $(u_1,u_2)\in\mathcal A$, set
$(n,w)=(u_1,u_2)2^kM_0L^N$.  If a bounded local observable $F$ can be
placed in a rectangle of time extent at most $n/2$ and spatial extent
at most $w/2$, then
\begin{equation}
 |\mu_{\beta_N;n,w}(F)-\omega_{\beta_N}(F)|
       \le C\norm{F}_\infty e^{-c2^k}.
 \label{eq:trace-uniform-volume}
\end{equation}
The same estimate holds for a torus obtained by closing time with any
integer spatial translation, provided the support has a lift of time
extent at most $n/2$ and spatial extent at most $w/2$.

For bounded local observables $F,G$ on the plane whose supporting time
slabs are separated by $d\ge0$ transfer edges,
\begin{equation}
 |\Cov_{\omega_{\beta_N}}(F,G)|
    \le\norm{F}_\infty\norm{G}_\infty
       e^{-c d/(M_0L^N)}.
 \label{eq:orig-26}
\end{equation}
The constant in this last estimate may be taken to be $c=\log2$.
\end{theorem}

\begin{proof}
First compare $(n,w)$ with $(2n,w)$.  Both torus expectations are
within $2\norm{F}_\infty s_{n/2}(w)$ of the cylinder expectation at
circumference $w$, by \eqref{eq:trace-torus-cylinder}.  The auxiliary
aspect $(u_1/2,u_2)$ and Proposition~\ref{prop:trace-doubling} give
\[
 s_{n/2}(w)
    \le e^{p(n/2,w)}-1
    \le e^{\Delta_{\beta_N}(n/2,w)}-1
    \le C e^{-c2^k}.
\]
Next compare $(2n,w)$ with $(2n,2w)$, applying the same argument in
the spatial direction.  The needed excited trace has running length
$w/2$ and transverse circumference $2n$; it is bounded using
$q(2n,w/2)\le\Delta_{\beta_N}(2n,w/2)$ and the auxiliary aspect
$(2u_1,u_2/2)$.  Thus consecutive simultaneous doublings differ by
at most $C\norm{F}_\infty e^{-c2^k}$.  Sum this estimate over all
subsequent doublings.  At fixed $N$, the limit is the periodic
infinite-plane state from Section~\ref{sec:preliminary}, giving
\eqref{eq:trace-uniform-volume}.

For completeness, this identification also holds for bounded
measurable local observables.  At fixed cutoff and fixed finite
support, nearest-neighbour conditional densities give a uniform upper
bound, independent of the enclosing torus, for the marginal density
relative to product area measure.  Approximation in that reference
measure extends the identity from continuous to bounded measurable
observables.  The direct transfer estimates above already have the
same supremum-norm bound for such observables.

Now close time with a spatial shift $s\in\mathbb Z/w\mathbb Z$.
Let $U_s$ be its unitary action on $L^2((S^2)^w)$.  Translation
invariance gives $U_sT_w=T_wU_s$.  Simplicity and positivity of the
ground state give $U_s\Omega_w=\Omega_w$.  Choosing the closing seam
outside the supporting slab, the normalized partition function and
the inserted numerator are
\[
 \Tr(T_w^nU_s),\qquad \Tr(A_FT_w^{n-r}U_s).
\]
Their ground-state contributions are $1$ and $\omega_w(F)$,
respectively, and their remaining contributions have absolute values
at most $s_n(w)$ and $\norm{F}_\infty s_{n-r}(w)$.  The half-period
bound just proved ensures $s_n(w)<1/2$.  Consequently
\[
 |\mu_{\beta_N;n,w}^{(s)}(F)-\omega_w(F)|
     \le4\norm{F}_\infty s_{n/2}(w).
\]
Comparison with the unshifted torus, followed by
\eqref{eq:trace-uniform-volume}, proves the assertion for shifted
closing conditions, uniformly in $s$.

Finally put $n_0=M_0L^N$.  The square reference-box test and
\eqref{eq:trace-width-gap} imply, at every circumference
$w=2^kn_0$,
\[
 \norm{T_w^d-P_w}\le e^{-(\log2)d/n_0}.
\]
For two ordered slab insertions on the cylinder,
\[
 \Cov_{\omega_w}(F,G)
   =\inner{\Omega_w}{A_{\overline F}(T_w^d-P_w)A_G\Omega_w}.
\]
The insertion norm estimate therefore proves
\eqref{eq:orig-26} on these cylinders with $c=\log2$.
Equation~\eqref{eq:trace-torus-cylinder}, using the half-period
trace bound, identifies their local limits with the square-torus
limit.  Passage to that limit proves the assertion on the plane.
\end{proof}

\subsection{A uniform physical gap}

The slab estimate controls the whole transfer spectrum, because it
applies to all bounded local observables.  Its spectral consequence is
therefore stronger than a bound on the spin two-point function alone.

\begin{corollary}[A uniform physical gap]
\label{cor:trace-physical-gap}
For $N\ge K_0$, let $\mathcal T_N$ be the one-step transfer contraction
in the reflection-positive Hilbert space of $\omega_{\beta_N}$, and
let $\Omega_N$ be its vacuum vector.  Then
\[
 \spec\bigl(\mathcal T_N|_{\Omega_N^\perp}\bigr)
    \subset\left[0,e^{-(\log2)/(M_0L^N)}\right].
\]
At lattice spacing $a_N=L^{-N}$, the lower energy scale in physical
units is consequently
\begin{equation}
 \frac{1}{a_N}\frac{\log2}{M_0L^N}
       =\frac{\log2}{M_0}>0.
 \label{eq:trace-physical-scale}
\end{equation}
\end{corollary}

\begin{proof}
Link reflection positivity follows from positive semidefiniteness of
the crossing-bond kernel $e^{\beta s\cdot s'}$; site reflection
positivity follows by conditional factorization across the fixed row.
Both pass to the periodic plane limit and make $\mathcal T_N$ a
positive semidefinite self-adjoint contraction.
A centered bounded local observable creates a vector whose transfer
autocorrelation is the covariance of two reflected slabs.
Equation~\eqref{eq:orig-26} bounds this autocorrelation by
$C_\psi e^{-(\log2)j/(M_0L^N)}$, with $C_\psi$ depending on the vector
and its slab thickness.  The spectral measure is positive.  Positive
mass on any interval $[r,1]$ with
$r>e^{-(\log2)/(M_0L^N)}$ would give a lower bound proportional to
$r^j$, contradicting that decay.  Centered bounded local vectors are
dense in the vacuum orthogonal subspace.  Boundedness of spectral
projections therefore gives the asserted exclusion for the whole
subspace, as in \Rref{prop:criterion}.  Dividing the negative logarithm
of a nonzero transfer spectral value by $a_N$ gives the physical
energy bound \eqref{eq:trace-physical-scale}.
\end{proof}

The constants and the physical volume threshold have been fixed before
removing the cutoff.  Continuum convergence is established in
Section~\ref{sec:continuum}; Section~\ref{sec:os} then transfers this
cutoff-independent positive scale to the continuum Hamiltonian.

\subsection{Tilted periods}

The only nonrectangular periods needed in the construction are covered
by the shifted closing condition in Theorem~\ref{thm:trace-uniform}.
Let $M=2^kM_0$.  A square of physical side $5M$ tilted by $O_+$ has
period vectors
\[
 a=M(3,4),\qquad b=M(-4,3)
\]
in the regulator axes, with coordinate pairs in this paragraph written
as $(\mathrm{space},\mathrm{time})$.  The integer change of basis
\[
 3a-4b=(25M,0),\qquad a-b=(7M,M)
\]
has determinant one.  The same torus therefore has spatial
circumference $25M$, time height $M$, and a spatial shift $7M$ upon
closing time.  Dividing by $a_N=L^{-N}$ makes all periods and shifts
integers.  Reflected versions, including $O_-$ in the other regulator
frame, have the same circumference and height, with a possibly
different shift.  A fixed compact physical support has a lift satisfying
both extent bounds for all sufficiently large $k$, uniformly in $N$.
Thus the finite list containing
$(1,25)$ and its auxiliary aspects supplies all the tilted-torus
comparisons needed later.
